%%%%%%%%%%%%%%%%%%%%%%%%%%%%%%%%%
%
%       EXERCÍCIO
%
%%%%%%%%%%%%%%%%%%%%%%%%%%%%%%%%%

\definevalorquestao

\begin{exercicioBanco}[\valorquestao]
    O número de divórcios por indivíduo adulto casado, em certa comunidade, foi modelado pela variável aleatória \(D\), cuja função de probabilidade é apresentada a seguir:

\[
\begin{array}{c|cccc}
D   & 0 & 1 & 2 & 3 \\ \hline
p_i & 0.5 & 0.4 & 0.05 & 0.05
\end{array}
\]

Uma amostra, representada por \((D_1, D_2)\), foi sorteada com dois desses indivíduos e os seguintes estimadores, para a média de divórcios, foram considerados:
\[
\hat{\mu}_1 = \sqrt{D_1D_2}
\quad \text{e} \quad
\hat{\mu}_2 = \max(D_1, D_2) - \min(D_1, D_2).
\]

Para cada estimador, obtenha sua distribuição de probabilidade e verifique se é viciado.
\end{exercicioBanco}

